\section{Discussion}\label{sec:discussion}

For an engineer who sizes the interval of a soft sensor whose labels arrive late, the results
reduce to five rules.
\begin{enumerate}\setlength{\itemsep}{0pt}
\item Do not tune the step on history that ignores the delay, or by the rule for uncorrelated
residuals. On the simulated plant they left 6.9 and 18.2\,\% of the intervals uninformative
against a 2\,\% budget, at coverages of 0.897 and 0.899, so report $U$ beside the coverage.
\item Tune the step by replaying the calibration residuals through the delayed loop at the plant's
dead time; the step it chooses lies far below the stability limit.
\item Compensate the predictable part of the residual first: at the same budget the interval was
18\,\% narrower, its misses less persistent, and a larger step held the budget.
\item Before choosing or moving an analyzer, read the width that compensation would buy off the
autocorrelation of the calibration residuals (Eq.~\eqref{eq:p3}); on average, the gain was never
smaller than predicted.
\item Where the evaluation data resemble the calibration data, compensated split conformal may be
enough. Where they seem not to, as on the debutanizer stream (split conformal 0.750), the loop raised
the coverage to 0.865 and, compensated, to 0.900, against 0.876 for compensated split conformal
(Fig.~\ref{fig:debutanizer}); the 95\,\% block-bootstrap intervals of both gains include 0.
\end{enumerate}

\emph{Why the width prediction errs on the safe side.} Eq.~\eqref{eq:p3} assumes that the
compensated and uncompensated scores share one standardised shape. With other shapes the width
ratio can differ from the ratio of standard deviations: for AR(1) residuals with $\phi = 0.9$ and
uniform or $t_5$ innovations, simulated over 2 million steps before the study, the width ratio at
one step ahead lay 4.2 to 5.5\,\% below it, the same sign and size as the errors at the longer
label ages here (0.2 to 5.4\,\%). The larger errors (18 to 44\,\%) suggest two further sources,
neither tested here. At the shortest label ages on the Tennessee Eastman process the analyzer holds
each reading, with its noise, for 0.25~h, so the residual one update ahead is mostly predictable and
its prediction error far from Gaussian, a stronger shape effect than the simulated one. On the
debutanizer the prediction uses the calibration residuals, and a slow shift in the evaluation rows
widens the uncompensated interval, which must cover it, more than the compensated one, whose
predictor follows it; the split-conformal coverage of 0.750 points to such a shift.

\emph{Limits.} The diagnosis rests on a simulator: the Tennessee Eastman runs share one plant model
and three operating points, come from the streams that did not shut down (36 of 76 candidates),
and share a tuned step within each operating point. The debutanizer is one stream, on which F2 had a
power of about 0.1. The dead times follow from documented delays through a timing convention
(Table~\ref{tab:delays}), which a plant must check against its own data historian. The ridge
penalty sat at the top of its grid in 10 of the 15 runs, so the point model may be
over-regularised; the gradient-boosting arm gave the same results on every rule computed for it.
The baselines were configured from their papers and code before the data, and other settings may
serve them better. Random and irregular
delays \citep{shardt2016,guo2014} were not studied.

